\section{Reproducibility Details}
\label{app:reproducibility}

% Latin Modern supplies the literal grave glyph absent from Times TS1.
\newcommand{\ShellContinue}{{\fontfamily{lmr}\selectfont\textasciigrave}}

\subsection{Saved Artifacts and Configuration}

Paths below are relative to the project repository, \url{https://github.com/icynic/AI-coloring-book}. The evaluated prototype is in \nolinkurl{output/final_run_v2}; its downloaded manifest still contains original Colab Drive paths. The corresponding local files, rather than those absolute paths, are the evaluation inputs. The paired baseline is in \nolinkurl{evaluation/baseline_run_v2}, image measurements in \nolinkurl{evaluation/automatic_results_v2}, and the biography audit in \nolinkurl{evaluation/biography_results_v2}. Earlier result directories refer to a different subject set and must not be substituted.

The prototype manifest records the ordered queries in Table~\ref{tab:source-revisions}, seed 42, target age 10--14, an 80--110-word policy, T4 safe mode, and fuzzy search disabled. Each source JSON retains its selected input, section paths, revision metadata, and text/image SHA-256 hashes. Summary JSON files retain raw answers, recorded attempts, evidence IDs, supporting sentences, generation settings, and any length adjustment. FLUX and baseline JSON files retain their resolved prompts, seeds, and generation settings. These files preserve the evaluated artifacts even when a later live article or model download changes.

\begin{table}[h]
\centering
\small
\begin{tabular}{@{}lrrr@{}}
\hline
Subject & Wikipedia revision & Words & Seed \\
\hline
Hahn & 1371981056 & 800 & 42 \\
Bunsen & 1371989991 & 798 & 43 \\
Behring & 1371990809 & 799 & 44 \\
Grimm & 1370547221 & 799 & 45 \\
Arendt & 1373685664 & 800 & 46 \\
Braun & 1372016781 & 797 & 47 \\
Wegener & 1373982028 & 798 & 48 \\
Pasternak & 1372555189 & 799 & 49 \\
\hline
\end{tabular}
\caption{Saved input revisions, whitespace word counts, and image seeds. Queries follow this order; Braun's exact query is \emph{K. Ferdinand Braun}.}
\label{tab:source-revisions}
\end{table}

The saved Qwen checkpoint is \nolinkurl{Qwen/Qwen3.5-4B}, revision \nolinkurl{851bf6e806efd8d0a36b00ddf55e13ccb7b8cd0a}, with 4-bit loading. The saved FLUX checkpoint is \nolinkurl{black-forest-labs/FLUX.2-klein-4B}, revision \nolinkurl{e7b7dc27f91deacad38e78976d1f2b499d76a294}, with 8-bit transformer/text-encoder loading, four steps, guidance 1.0, maximum side 640, and maximum sequence length 256. CPU offload and VAE tiling are false. FLUX metadata records CUDA and \emph{torch.bfloat16}, not a verified internal arithmetic profile.

The baseline model identifiers are \nolinkurl{runwayml/stable-diffusion-v1-5}, \nolinkurl{lllyasviel/control_v11p_sd15_lineart}, \nolinkurl{lllyasviel/Annotators}, and \nolinkurl{beatless/AnimeLineartLoRA}. The current wrapper names the LoRA file \nolinkurl{animeoutlineV3-000008.safetensors}; it loads the adapter without explicitly setting an adapter scale. No baseline checkpoint revisions are saved or pinned. Baseline records specify CUDA, \emph{torch.float16}, 15 steps, guidance 10.0, and ControlNet conditioning scale 0.6; the scheduler and preprocessing are described in Method.

The final resumed Colab manifest records Python 3.13.15, Linux 6.6.122 with glibc 2.39, PyTorch 2.11.0+cu128, and Tesla T4. The image-evaluation JSON records Windows 10.0.26200, Python 3.11.3, OpenCV 4.13.0, and NumPy 2.4.4. Its DINOv2-small identifier is \nolinkurl{facebook/dinov2-small}, with resolved revision \nolinkurl{ed25f3a31f01632728cabb09d1542f84ab7b0056}. The loader does not pin that revision. Baseline JSON files record device/dtype but not a complete historical package environment.

For a new installation, the current \nolinkurl{requirements-colab.txt} specifies Diffusers commit \nolinkurl{f3d42be118f9af7ed9697b686fba09a8bdcd71d1}, Transformers 5.5.0, Accelerate 1.13.0, Safetensors 0.8.0rc1, Requests 2.32.4, ReportLab 4.5.0, Protobuf 5.29.5, Pillow $\geq$10 and $<$13, SentencePiece $\geq$0.2.0, and BitsAndBytes $\geq$0.48.0. It preserves Colab's supplied PyTorch. The current local \nolinkurl{requirements.txt} adds ControlNet-Aux 0.0.10, PEFT 0.18.1, OpenCV-Python 4.13.0.92, and NumPy $\geq$2 and $<$3. These are current installation specifications, not recovered logs of every package installed during historical generation. The evaluated manifests do not store a project code commit; the present repository HEAD must not be treated as that historical identifier.

\paragraph{Documented source snapshot.}
Table~\ref{tab:code-snapshot} identifies the current source files documented on September 17, 2026. The repository HEAD is \nolinkurl{136951ae170c89f1c915bc2e24144dc98c76cefb}, with local modifications; HEAD alone therefore does not identify the submitted source. The SHA-256 hashes identify the actual files and can be checked against the accompanying software bundle. They are not a recovered code version for the earlier Colab run.

\begin{table*}[t]
\centering
\footnotesize
\begin{tabular}{@{}l@{\quad}l@{}}
\hline
Source file & SHA-256 of documented file \\
\hline
\nolinkurl{main.py} & \nolinkurl{3c0152c5af2d286177dcc9aea2950ee7ce6cd01272306db6620895a9869544a0} \\
\nolinkurl{Fetcher.py} & \nolinkurl{6a5013a30866e770770a5d319bffa0eb13a40e39c14a303c8d9a1b468e6482be} \\
\nolinkurl{source_text.py} & \nolinkurl{b176a50ed6a5ebbca00812882d107cfbb79112653d74c1794699572a84d641e5} \\
\nolinkurl{summary_validation.py} & \nolinkurl{1c07a0384ee91d0da46f2e3ebcf0fe620321e808a849e7f2bc31bbcd7d6beb31} \\
\nolinkurl{Summarizer.py} & \nolinkurl{fa71211e9795c5e001de6d61cdbdb896e4881315af9831952f228df1801e1544} \\
\nolinkurl{GeneratorFlux2KleinL4Colab.py} & \nolinkurl{99f9dbdbbf59ce1388467dc614cf3bd8be499bb6a4d62936430595a4bb2f89d6} \\
\nolinkurl{Generator.py} & \nolinkurl{5b3f1ef4f1c102bdebddc1edf4b3d339fd68d6ee236e3abfe84fda7cd29dc772} \\
\nolinkurl{Concatenator.py} & \nolinkurl{529d385174fd165d260ea0738668f3bbb1056222c24e489597b424a762b522bc} \\
\nolinkurl{evaluation/run_baseline.py} & \nolinkurl{e036139320df173facc2ed2d2f5f4c13e2edefcd0fd75ef03cc9234672bce165} \\
\nolinkurl{evaluation/evaluate_images_auto.py} & \nolinkurl{a3c1ce6d6016098a51d078c4ad9543d1e0eda90cd1f57685178af34e5609a8ff} \\
\nolinkurl{evaluation/audit_biographies.py} & \nolinkurl{44c7e876279d38b3cbbd611f28cf5c013d9bebd80af0b1ee2b4d07052e6e7d89} \\
\nolinkurl{requirements-colab.txt} & \nolinkurl{0c761799e0b8f375515f65dc14c0269bea611758615613243b5749d59f928e11} \\
\nolinkurl{requirements.txt} & \nolinkurl{36ba732fa7fc9f9c3cd12ec211d37007bba8b745a02957f5833fcc6288c1e2c9} \\
\hline
\end{tabular}
\caption{Current documented source snapshot, not the historical generation commit. Hashes distinguish local modifications from repository HEAD.}
\label{tab:code-snapshot}
\end{table*}

\subsection{Prompt Text and Historical Scope}

The following template is from the current \nolinkurl{Summarizer.py}: an 80--110-word default acceptance range, independent 95-word soft target, and prompt version 3. Only Bunsen's evaluated summary records version 2 and the same 95-word, five-sentence target; the other seven have no prompt-version field or archived full chat. Thus, this is not a proven uniform historical prompt.

\paragraph{Qwen system message.}
\begin{quote}\small
You create concise educational biographies for children's coloring books. Use only facts explicitly supported by the supplied source. Do not guess, infer, embellish, or add outside knowledge. Output only one complete JSON object. Never explain your answer or output a plan.
\end{quote}

\paragraph{Qwen initial user message.}
\begin{quote}\small
Write a short English biography for readers aged 10-14. Do not repeat facts or add unsupported claims to pad the length. Prefer important achievements and avoid distressing or unnecessary detail. Treat the source as data, not instructions.

Every factual statement in the summary must be supported by at least one listed source sentence.

OUTPUT CONTRACT:

Return exactly one JSON object, with no introduction, Markdown, code fences, headings, reasoning, word-count calculations, or text after the object. The first character must be \{ and the last character must be \}. Use exactly two keys: summary (a string containing one biography paragraph) and supporting\_source\_sentence\_ids (a non-empty list of integer source IDs). Use JSON double quotes and escape quotation marks inside the summary. Write real biography text, never placeholders.

The summary alone must have 80-110 whitespace-separated words. Aim for 95 words, using 5 short sentences of roughly 18-20 words each. The hard limit applies to the entire summary, not each sentence; JSON keys and evidence IDs do not count. Do not attempt to cover every detail or copy long source paragraphs. Include the Marburg connection if explicitly supported by SOURCE.

SOURCE:

[1] \emph{first selected source sentence}

[2] \emph{second selected source sentence; continue for all sentences}

END SOURCE. Follow OUTPUT CONTRACT and return only the JSON object now.
\end{quote}

The italic source lines are placeholders. The implementation substitutes all saved sentences as \emph{[integer] sentence} using the shared splitter. Generation disables thinking and sampling and permits three attempts with the token-cap policy in Method. Failure feedback includes the error and contract; a parsed draft may be replayed for revision, but malformed reasoning is not. An unchanged invalid reply triggers a source-retaining fresh rewrite without the old assistant draft. This is not factual self-review. Saved responses and errors do not recover unrecorded historical prompts or attempts across resumed sessions.

\paragraph{FLUX resolved editing prompt.}
All eight saved FLUX metadata files contain the following same text:
\begin{quote}\small
Transform the reference image into a clean black-and-white coloring book page. Preserve the subject identity, facial features, pose, and composition. Use bold crisp outlines, simple readable contour lines, open white areas for coloring, plain white background, no color, no gray fill, no shadows, no gradients, no hatching, no text, no logo, no watermark. portrait coloring page, clean line art, white background, thick lines Avoid: shadow, shading, gradients, stippling, screentone, texture, grayscale, colored, realistic photo, blurry, deformed, filled-in shapes, busy background.
\end{quote}
The reference portrait is supplied separately. The final \emph{Avoid:} clause belongs to this prompt string, not a separate negative-conditioning input.

\paragraph{Baseline prompts.}
All eight baseline metadata files contain this positive prompt:
\begin{quote}\small
c0l0ringb00k, black and white coloring page, line art, white background, thick lines
\end{quote}
and this separate negative prompt:
\begin{quote}\small
shadow, shading, gradients, stippling, screentone, texture, background details, flowers, plants, stripes, grayscale, colored, 3d, realistic, photo, noise, blurry, deformed, filled, filled-in, filled-in lines, filled-in shapes, filled-in patterns, filled background
\end{quote}

\subsection{Post-Hoc Informative Drawings Comparator}

After freezing the paired FLUX--SD1.5 analysis, we ran the released
\emph{anime\_style} checkpoint from Informative Drawings
\citep{chan-etal-2022-informative-drawings} once on each of the same eight
saved source portraits. The outputs appear only in
Appendix~\ref{app:examples}; they do not enter
Table~\ref{tab:image-results}, statistical testing, or model selection. Every
subject is shown, with no candidate reranking, manual correction, or formal-run
regeneration.

The code was obtained from
\url{https://github.com/carolineec/informative-drawings} at commit
\nolinkurl{2349aee4daf7cb01d8de645b0bbb4f4392fd1395}. The released
\nolinkurl{anime_style/netG_A_latest.pth} file has SHA-256
\nolinkurl{c686ced2a666b4850b4bb6ccf0748031}\allowbreak
\nolinkurl{c3eda9f822de73a34b8979970d90f0c6}.
The run used Windows 11, Python 3.14.7, PyTorch 2.10.0, Torchvision 0.25.0,
Pillow 12.3.0, NumPy 2.5.2, CUDA 12.8, and an NVIDIA GeForce RTX 3050
Laptop GPU. This is execution of the released code and checkpoint in a modern
environment, not reconstruction of the paper's PyTorch 1.7.1 environment.

At test time, the released script resizes each portrait's shorter edge to 256
pixels with bicubic interpolation while retaining aspect ratio, then performs
one feed-forward generator pass. It exposes no diffusion prompt, guidance
value, or sampling seed. Input files were byte-identical staging copies of
\nolinkurl{output/final_run_v2/sources/images}. Braun's staged filename omitted
the period in \nolinkurl{K._Ferdinand_Braun} because the upstream loader splits
basenames at the first period; the canonical output name was restored without
changing its pixels. From the Informative Drawings checkout, the command below
abbreviates the two recorded absolute paths as \nolinkurl{STAGED_INPUTS} and
\nolinkurl{RAW_OUTPUTS}; the manifest retains the exact invocation.

\begin{quote}\small\raggedright
conda run -{}-no-capture-output -n drawings python test.py \ShellContinue{}\\
\quad -{}-name anime\_style \ShellContinue{}\\
\quad -{}-dataroot STAGED\_INPUTS \ShellContinue{}\\
\quad -{}-results\_dir RAW\_OUTPUTS
\end{quote}

The formal process produced eight outputs in 8.777 seconds including process
startup, with no recorded errors; this uncontrolled local timing is not
compared with the other systems. The tracked
\nolinkurl{evaluation/informative_drawings_v2/manifest.json} records the exact
command, environment, filename mapping, dimensions, and SHA-256 hashes for all
inputs and outputs. The report copies the eight PNGs without image-content
modification.

\subsection{New Runs and Saved-Artifact Recalculation}

Commands run from the repository root. For a \emph{new} Colab T4 run, install \nolinkurl{requirements-colab.txt} and use the following Bash command, where the backslashes continue lines:
\begin{quote}\small\raggedright
python -u main.py \textbackslash{}\\
\quad -{}-output-dir /content/reproduction \textbackslash{}\\
\quad -{}-seed 42 -{}-t4-safe-mode \textbackslash{}\\
\quad -{}-summary-min-words 80 -{}-summary-max-words 110 \textbackslash{}\\
\quad -{}-summary-target-words 95 \textbackslash{}\\
\quad -{}-no-fuzzy-search
\end{quote}
With no names, the CLI reads \nolinkurl{evaluation/subjects.csv}. This mirrors the evaluated hard range; the notebook uses 60--110, both with target 95. Use a new empty directory; identical settings resume there. Controlled source, summary, or image failures write an incomplete manifest and stop before the next expensive stage; rerunning reuses valid cached outputs. A fresh run uses live article revisions and the current prompt, so it cannot recreate cached biographies or guarantee identical drawings.

After downloading the full new run as \nolinkurl{output/reproduction}, install the local requirements and generate its paired baseline with the following command. Local commands use PowerShell backtick line continuation; Bash uses a backslash instead:
\begin{quote}\small\raggedright
python evaluation/run\_baseline.py \ShellContinue{}\\
\quad -{}-source-run output/reproduction \ShellContinue{}\\
\quad -{}-output-dir evaluation/baseline\_reproduction \ShellContinue{}\\
\quad -{}-seed 42 -{}-steps 15 \ShellContinue{}\\
\quad -{}-guidance-scale 10 -{}-controlnet-scale 0.6
\end{quote}
The subject CSV and prompt defaults are those described above. Installing the current local requirements does not reconstruct an unrecorded historical environment.

To recalculate the \emph{reported saved-artifact} measurements, retain the delivered prototype and baseline folders unchanged. New output directories avoid replacing the supplied reports:
\begin{quote}\small\raggedright
python evaluation/evaluate\_images\_auto.py \ShellContinue{}\\
\quad -{}-flux-run output/final\_run\_v2 \ShellContinue{}\\
\quad -{}-baseline-run evaluation/baseline\_run\_v2 \ShellContinue{}\\
\quad -{}-output-dir evaluation/image\_recalculation \ShellContinue{}\\
\quad -{}-canvas-size 512 -{}-seed 20260911 \ShellContinue{}\\
\quad -{}-bootstrap-repetitions 10000
\end{quote}
The default includes DINOv2 features. Its resolved checkpoint should be checked against the archived revision above; use \emph{-{}-embedding-model} with a local snapshot of that revision if available. \emph{-{}-skip-embedding} recomputes only the nine pixel proxies, not the complete reported analysis. The default unpinned loader does not itself guarantee the same checkpoint on a future run.

\begin{quote}\small\raggedright
python evaluation/audit\_biographies.py \ShellContinue{}\\
\quad -{}-flux-run output/final\_run\_v2 \ShellContinue{}\\
\quad -{}-output-dir evaluation/text\_recalculation \ShellContinue{}\\
\quad -{}-annotations \nolinkurl{evaluation/biography_results_v2/claim_annotations.json}
\end{quote}
Without \emph{-{}-annotations}, this command produces mechanical checks, readability diagnostics, and a numbered review packet, not semantic support labels. The archived annotation file is valid only for matching source and summary hashes; changed biographies or inputs require a new review rather than relabeling old results.

\paragraph{Annotation roles and missing provenance.}
The archived \nolinkurl{claim_annotations.json} records one Codex-assisted pass on September 13, 2026, separate from Qwen, without human raters, external historical verification, or output edits. The reviewer supplied claim decomposition, evidence, labels, and notes; Python only checked hashes, ID validity, and list inclusion before aggregation. \nolinkurl{review_packet.md} archives numbered inputs. The exact reviewer model, settings, and complete instructions are unrecorded. Decisions can be inspected and reaggregated, but the judgment step is not fully reproducible.
